% GENERATED FILE. Edit canonical lessons, then run npm run generate:books.
% canonical-source-hash: fnv1a64:fd77a039e928a933
% canonical-lessons: MW-C211-map, MW-C211-napno, MW-C211-moklo, MW-C211-tang, MW-C211-samavno

\chapter{Roomy or Tight}
\label{ch:mw-roomy-or-tight}
\hlchaptermodality{211}

\begin{chapteropening}
\textbf{By the end of this chapter:} \emph{I can say whether clothes fit.}

Complete the last lesson of chapter 211: I can say whether clothes fit.
\end{chapteropening}

\section[\texorpdfstring{\mw{माप} (māp)}{माप (māp)}]{\mw{माप} (māp) — a size; a measure}

\label{lesson:MW-C211-map}



\begin{quote}
Before the new one: say the Marwadi for fake, then the Marwadi for to pick out.
\end{quote}

\subsection*{You'll want to know: \mw{माप}}
\textbf{\mw{माप}} — \emph{māp} — \textquotedblleft{}a size; a measure\textquotedblright{}.

\begin{quote}
\mw{ये} \mw{जूती} \mw{म्हारो} \mw{माप} \mw{कोनी।}
\end{quote}

\begin{quote}
\emph{ye jūtī mhāro māp konī.}
\end{quote}

\begin{quote}
These shoes aren't my size.
\end{quote}

A masculine noun.

\begin{cousinweb}
\textbf{\mw{माप}} (\emph{māp}) is inherited from Sanskrit \emph{māpya}, \textquotedblleft{}to be measured\textquotedblright{}. CDIAL gives no Marwari form in this entry, but lists Hindi \emph{māp}, \textquotedblleft{}measurement, measure\textquotedblright{} (CDIAL 10053).
\end{cousinweb}

\begin{grammarlens}[title={What you've built}]
One more word for this chapter.
\end{grammarlens}

\subsection*{Your turn}
\begin{itemize}
\raggedright
  \item \emph{Say it:} \emph{māp}
  \item \emph{Say it:} \emph{māp}, once more
  \item \emph{Say it:} \emph{ye jūtī mhāro māp konī.} — These shoes aren't my size.
  \item \emph{Recall:} say the Marwadi for fake, then the Marwadi for to pick out, then say \emph{māp} again
\end{itemize}

\begin{tcolorbox}[breakable,colback=teal!4,colframe=teal!35!black,title={Before you move on}]
What does \mw{माप} mean? (\textquotedblleft{}A size; a measure\textquotedblright{}.) Say it once more.
\end{tcolorbox}
\section[\texorpdfstring{\mw{नापणो} (nāpṇo)}{नापणो (nāpṇo)}]{\mw{नापणो} (nāpṇo) — to measure}

\label{lesson:MW-C211-napno}



\begin{quote}
Before the new one: say the Marwadi for to pick out, then the Marwadi for a size.
\end{quote}

\subsection*{You'll want to know: \mw{नापणो}}
\textbf{\mw{नापणो}} — \emph{nāpṇo} — \textquotedblleft{}to measure\textquotedblright{}.

\begin{quote}
\mw{मिस्त्री} \mw{भींत} \mw{नापै} \mw{है।}
\end{quote}

\begin{quote}
\emph{mistrī bhī̃t nāpai hai.}
\end{quote}

\begin{quote}
The mason measures the wall.
\end{quote}

Past (I …): \mw{म्हैं} \mw{कमरो} \mw{नाप्यो}, \emph{mhaĩ kamro nāpyo}.

\begin{cousinweb}
\textbf{\mw{नापणो}} (\emph{nāpṇo}) is inherited from Sanskrit \emph{jñāpyate}, \textquotedblleft{}is made known\textquotedblright{}. CDIAL gives no Marwari form in this entry, but lists Hindi \emph{nāpnā}, \textquotedblleft{}to measure\textquotedblright{} (CDIAL 5283).
\end{cousinweb}

\begin{grammarlens}[title={What you've built}]
One more word for this chapter.
\end{grammarlens}

\subsection*{Your turn}
\begin{itemize}
\raggedright
  \item \emph{Say it:} \emph{nāpṇo}
  \item \emph{Say it:} \emph{nāpṇo}, once more
  \item \emph{Say it:} \emph{mistrī bhī̃t nāpai hai.} — The mason measures the wall.
  \item \emph{Recall:} say the Marwadi for to pick out, then the Marwadi for a size, then say \emph{nāpṇo} again
\end{itemize}

\begin{tcolorbox}[breakable,colback=teal!4,colframe=teal!35!black,title={Before you move on}]
What does \mw{नापणो} mean? (\textquotedblleft{}To measure\textquotedblright{}.) Say it once more.
\end{tcolorbox}
\section[\texorpdfstring{\mw{मोकळो} (mokḷo)}{मोकळो (mokḷo)}]{\mw{मोकळो} (mokḷo) — roomy; loose}

\label{lesson:MW-C211-moklo}



\begin{quote}
Before the new one: say the Marwadi for a size, then the Marwadi for to measure.
\end{quote}

\subsection*{You'll want to know: \mw{मोकळो}}
\textbf{\mw{मोकळो}} — \emph{mokḷo} — \textquotedblleft{}roomy; loose\textquotedblright{}.

\begin{quote}
\mw{ये} \mw{कुरतो} \mw{मोकळो} \mw{है।}
\end{quote}

\begin{quote}
\emph{ye kurto mokḷo hai.}
\end{quote}

\begin{quote}
This kurta is roomy.
\end{quote}

\begin{cousinweb}
\textbf{\mw{मोकळो}} (\emph{mokḷo}) is inherited from a reconstructed \emph{*mukna}, \textquotedblleft{}loosed\textquotedblright{}; CDIAL lists Old Marwari \emph{mokaï}, \textquotedblleft{}sends\textquotedblright{}, and Hindi \emph{moklā}, \textquotedblleft{}loose, free\textquotedblright{}, under it (CDIAL 10157).
\end{cousinweb}

\begin{grammarlens}[title={What you've built}]
One more word for this chapter.
\end{grammarlens}

\subsection*{Your turn}
\begin{itemize}
\raggedright
  \item \emph{Say it:} \emph{mokḷo}
  \item \emph{Say it:} \emph{mokḷo}, once more
  \item \emph{Say it:} \emph{ye kurto mokḷo hai.} — This kurta is roomy.
  \item \emph{Recall:} say the Marwadi for a size, then the Marwadi for to measure, then say \emph{mokḷo} again
\end{itemize}

\begin{tcolorbox}[breakable,colback=teal!4,colframe=teal!35!black,title={Before you move on}]
What does \mw{मोकळो} mean? (\textquotedblleft{}Roomy; loose\textquotedblright{}.) Say it once more.
\end{tcolorbox}
\section[\texorpdfstring{\mw{तंग} (taṅg)}{तंग (taṅg)}]{\mw{तंग} (taṅg) — tight}

\label{lesson:MW-C211-tang}



\begin{quote}
Before the new one: say the Marwadi for to measure, then the Marwadi for roomy.
\end{quote}

\subsection*{You'll want to know: \mw{तंग}}
\textbf{\mw{तंग}} — \emph{taṅg} — \textquotedblleft{}tight\textquotedblright{}.

\begin{quote}
\mw{ये} \mw{पजामो} \mw{तंग} \mw{है।}
\end{quote}

\begin{quote}
\emph{ye pajāmo taṅg hai.}
\end{quote}

\begin{quote}
These trousers are tight.
\end{quote}

\begin{cousinweb}
\textbf{\mw{तंग}} (\emph{taṅg}) is usually traced to Persian \emph{tang}, \textquotedblleft{}narrow, tight\textquotedblright{}.
\end{cousinweb}

\begin{grammarlens}[title={What you've built}]
One more word for this chapter.
\end{grammarlens}

\subsection*{Your turn}
\begin{itemize}
\raggedright
  \item \emph{Say it:} \emph{taṅg}
  \item \emph{Say it:} \emph{taṅg}, once more
  \item \emph{Say it:} \emph{ye pajāmo taṅg hai.} — These trousers are tight.
  \item \emph{Recall:} say the Marwadi for to measure, then the Marwadi for roomy, then say \emph{taṅg} again
\end{itemize}

\begin{tcolorbox}[breakable,colback=teal!4,colframe=teal!35!black,title={Before you move on}]
What does \mw{तंग} mean? (\textquotedblleft{}Tight\textquotedblright{}.) Say it once more.
\end{tcolorbox}
\section[\texorpdfstring{\mw{समावणो} (samāvṇo)}{समावणो (samāvṇo)}]{\mw{समावणो} (samāvṇo) — to fit (into)}

\label{lesson:MW-C211-samavno}



\begin{quote}
Before the new one: say the Marwadi for roomy, then the Marwadi for tight.
\end{quote}

\subsection*{You'll want to know: \mw{समावणो}}
\textbf{\mw{समावणो}} — \emph{samāvṇo} — \textquotedblleft{}to fit (into)\textquotedblright{}.

\begin{quote}
\mw{म्हारो} \mw{पग} \mw{जूती} \mw{मांय} \mw{कोनी} \mw{समावै।}
\end{quote}

\begin{quote}
\emph{mhāro pag jūtī mā̃y konī samāvai.}
\end{quote}

\begin{quote}
My foot won't fit in the shoe.
\end{quote}

Past (it …): \mw{वो} \mw{सामान} \mw{पेटी} \mw{मांय} \mw{समायो}, \emph{vo sāmān peṭī mā̃y samāyo}.

\begin{cousinweb}
\textbf{\mw{समावणो}} (\emph{samāvṇo}) is inherited from Sanskrit \emph{sammāti}, \textquotedblleft{}is contained in\textquotedblright{}; CDIAL lists Old Marwari \emph{samāi}, \textquotedblleft{}fits into, is contained in\textquotedblright{} (CDIAL 12975).
\end{cousinweb}

\begin{grammarlens}[title={What you've built}]
One more word for this chapter.
\end{grammarlens}

\subsection*{Your turn}
\begin{itemize}
\raggedright
  \item \emph{Say it:} \emph{samāvṇo}
  \item \emph{Say it:} \emph{samāvṇo}, once more
  \item \emph{Say it:} \emph{mhāro pag jūtī mā̃y konī samāvai.} — My foot won't fit in the shoe.
  \item \emph{Recall:} say the Marwadi for roomy, then the Marwadi for tight, then say \emph{samāvṇo} again
\end{itemize}

\begin{tcolorbox}[breakable,colback=teal!4,colframe=teal!35!black,title={Before you move on}]
What does \mw{समावणो} mean? (\textquotedblleft{}To fit (into)\textquotedblright{}.) Say it once more.
\end{tcolorbox}
